%% introduction.tex -- submission version

\section{Introduction}
\label{sec:introduction}

LLM agents increasingly rely on self-healing mechanisms: after a failed tool
call, test, or generated program, an orchestrator captures diagnostic content
and asks the model to try again. This improves robustness, but it also creates
a trust boundary. The diagnostic payload may contain external content---for
example, a fetched page, a file, a test fixture, or tool output---that an
attacker can influence. If the recovery path re-injects that payload in a
recovery frame that asks the model to act on it, the failure channel becomes a
prompt-injection path distinct from the normal forward input path. We call this
class recovery-channel prompt injection (\rci{}).

The security question is therefore not only whether a model follows an
injected directive. It is whether the framework preserves the origin and
authority of content while transporting it from an execution failure into the
next LLM turn. A provenance marker can indicate that content is external, but
the marker remains a probabilistic control: it changes model context without
enforcing an authorization decision at the action sink. Recovery actions that
change roles, tools, or topology therefore require a separate structural
authorization check, which we formulate as a design implication.

\paragraph{Research questions.}
We study three empirical questions and one design implication:

\begin{description}[leftmargin=*]
    \item[\textbf{RQ1 (Source-level pattern).}] How do self-healing frameworks represent the provenance and authority of failure content on recovery paths?
    \item[\textbf{RQ2 (Exploitability).}] Can the resulting trust gap be exploited through native recovery/error flows, and how does it compare with matched benign recovery?
    \item[\textbf{RQ3 (Mitigation boundary).}] How do provenance, role, and recovery-position changes affect attack success, and what is the observed limit of a defense-aware recovery screen?
    \item[\textbf{Design implication.}] Which recovery actions require structural reauthorization in addition to content-level handling?
\end{description}

\paragraph{Study design.}
RQ1 uses systematic multiple-case source-code analysis with qualitative coding
over six open-source frameworks fixed at recorded revisions. RQ2 uses a
canonical Reflexion matrix and native AutoGen and Aider executions, each paired
with benign controls. RQ3 uses the native Reflexion TrustOrigin port as its
primary provenance contrast, with message-level role/position screens, a
five-class defense-aware screen, and a legitimate-recovery task screen retained
as bounded supplementary checks. The topology-mutation/\reauthgate{} material
is retained as a design implication; it is not presented as a fourth,
independently completed empirical evaluation.

\paragraph{Software-engineering positioning.}
We do not claim a new underlying language-model failure mode. The contribution
is the identification and empirical characterization of a distinct
software-framework transport path and trust boundary: externally influenced
failure content is captured, reframed, and re-injected by recovery
infrastructure. This is a software-engineering study because its unit of
analysis is the recovery implementation, its cases are fixed framework
revisions, and its outputs are maintainer-facing design and testing guidance.

\paragraph{Contributions.}
The study makes four bounded contributions. First, it provides a recovery-path
trust-boundary taxonomy and a reproducible source-audit ledger. Second, it
measures native-framework recovery-path behavior across Reflexion, AutoGen, and
Aider, with framework/model/payload stratification and matched benign controls.
Third, it quantifies the native provenance treatment and the limits of
supplementary defense screens without overstating small or underpowered
contrasts. Fourth, it specifies structural
reauthorization as a complementary, LLM-independent control for privileged
recovery side effects.

The canonical Reflexion matrix spans ASR \ReflexionMinASR{}--\ReflexionMaxASR{}
across five model aliases, while its \ReflexionBenignTrials{} matched benign
trials contain no classified attack. The provider catalog and local artifacts
are disclosed in the provenance manifest. Because the provider exposes aliases
but not an alias-to-checkpoint mapping, the study reports provider-catalog
provenance rather than immutable provider snapshots; all claims are restricted
to the recorded artifacts and protocol labels.
